% Transcription — fonds Alexandre Grothendieck, shelfmark 161-2, batch 2 (pages 21-40).
% Established from the facsimile archives/batches/161-2/batch-02.pdf (a mirror of
% https://grothendieck.umontpellier.fr/161-2.pdf), per the skill
% transcribe-grothendieck. Pages 25, 32 and 35 are blank — absent. The batch holds
% three distinct runs: the algebraic-theories draft continued from batch 1 (pp. 21-31),
% two loose sheets on Lawvere theories and monadicity (pp. 33-34), and a spiral
% notebook on the classifying topoi of ring axioms (pp. 36-40).
% Pass: Opus 5 (claude-opus-5), 2026-09-02 — first pass, unchecked against the pages by a human.
\documentclass[11pt,a4paper]{article}
% Le préambule commun à toutes les transcriptions du fonds Grothendieck.
%
% Il tient en une idée : **l'incertitude se compose, elle ne se résout pas.**
% Une transcription qui lisse un mot douteux a détruit la seule chose pour
% laquelle on la faisait — un lecteur doit pouvoir savoir, à chaque endroit,
% ce qui était sur le feuillet et ce qui a été ajouté. Les six macros qui
% suivent sont donc l'appareil critique, et non de la décoration.
%
% Les mêmes commandes sont reconnues par `scripts/render.mjs`, qui produit la
% vue de lecture du site. Une macro ajoutée ici doit l'être aussi là-bas,
% faute de quoi le rendu échouera bruyamment — ce qui est voulu : un rendu
% silencieusement faux est pire qu'un rendu absent.

\ProvidesPackage{grothendieck}[2026/08/08 Transcriptions du fonds Grothendieck]

\usepackage{amsmath,amssymb,amsthm}
\usepackage{mathrsfs}
\usepackage{stmaryrd}
% Les diagrammes commutatifs sont partout dans ce fonds : c'est la langue
% même de la géométrie algébrique de Grothendieck, et une transcription qui
% les rendrait en prose ne transcrirait rien.
\usepackage{tikz-cd}
% Français par défaut — c'est la langue des feuillets — mais l'anglais est
% chargé aussi : la traduction et le résumé sont des documents anglais, et la
% typographie française (espaces avant les deux-points) y serait une faute.
% Ces documents basculent par \selectlanguage{english} en tête de corps.
\usepackage[english,french]{babel}
\usepackage{fontspec}
\usepackage{geometry}
\geometry{a4paper,margin=25mm,marginparwidth=28mm}
\usepackage{marginnote}
\usepackage{xcolor}
% ulem, not soul. `soul`'s \st and \ul are built on a character-by-character
% scan that XeLaTeX's font machinery breaks: the document fails with an
% undefined control sequence rather than degrading. ulem does the same two jobs
% and is Unicode-engine safe. `normalem` stops it hijacking \emph, which the
% transcriptions use for Grothendieck's own emphasis.
\usepackage[normalem]{ulem}
\usepackage{hyperref}
\hypersetup{colorlinks=true,linkcolor=black,urlcolor=black,pdfborder={0 0 0}}

% --- Métadonnées du lot -----------------------------------------------------
% Reprises telles quelles par le script de rendu, qui les affiche en tête de
% la vue de lecture. Elles fixent ce que le fichier prétend couvrir : sans
% elles, une transcription flotte sans point d'attache dans un fonds de seize
% mille feuillets.

\newcommand{\folder}[1]{\gdef\grothFolder{#1}}
\newcommand{\batch}[1]{\gdef\grothBatch{#1}}
\newcommand{\leaves}[2]{\gdef\grothFirst{#1}\gdef\grothLast{#2}}
\newcommand{\foldertitle}[1]{\gdef\grothFoldertitle{#1}}
\gdef\grothFolder{?}\gdef\grothBatch{?}\gdef\grothFirst{?}\gdef\grothLast{?}\gdef\grothFoldertitle{}

% --- Le résumé --------------------------------------------------------------
%
% Il ouvre la lecture modernisée, sous le titre « Résumé », et il est composé
% en petit. Ce n'est pas un ornement : c'est le seul passage du document qui
% ne soit pas des mathématiques, et un lecteur doit pouvoir le reconnaître
% avant de l'avoir lu — puis le sauter s'il connaît déjà le sujet, ou s'y
% arrêter s'il ne le connaît pas. Le filet qui le suit marque la reprise du
% corps.

\newenvironment{resume}
  {\small\setlength{\parskip}{0.45em}}
  {\par\medskip\hrule\medskip}

% --- Mots-clés ---------------------------------------------------------------
%
% En anglais, en clôture du résumé de la lecture modernisée : le vocabulaire
% moderne sous lequel chercher ce que les feuillets construisent. Le manifeste
% du site (`npm run manifest`) les extrait pour étiqueter la cote entière —
% cette ligne est la source unique des tags, il n'y a pas de fichier à côté.
\newcommand{\keywords}[1]{\par\smallskip\noindent
  {\footnotesize\textsc{Keywords} --- \textit{#1}}\par}

% --- L'appareil critique ----------------------------------------------------

% Le numéro de feuillet, en marge. C'est l'ancre qui permet de régler un
% désaccord sur une formule en regardant *un* feuillet plutôt qu'une chemise
% de six cents. Le site s'en sert aussi pour tourner les pages du fac-similé
% quand on fait défiler la transcription.
\newcommand{\leaf}[1]{%
  \marginnote{\footnotesize\color{gray!70}\textsf{#1}}%
  \label{leaf:#1}\ignorespaces}

% Illisible : on ne devine pas. Un mot inventé qui se lit comme les autres est
% le pire résultat possible.
\newcommand{\ill}{\textcolor{red!60!black}{[\,\ldots\,]}}

% Lecture douteuse : le mot est donné, et signalé comme incertain.
\newcommand{\uncertain}[1]{\uline{#1}}

% Ajout de l'éditeur — une lettre manquante, un mot sous-entendu.
\newcommand{\add}[1]{\textcolor{blue!50!black}{[#1]}}

% Biffé par Grothendieck lui-même. Ce qu'il a rayé fait partie du document :
% une rature dit souvent où la pensée a changé de direction.
\newcommand{\struck}[1]{\sout{#1}}

% Note du transcripteur, sur l'état matériel du feuillet.
\newcommand{\note}[1]{\par\smallskip{\small\color{gray!60!black}\textit{[#1]}}\par\smallskip}

% Note marginale de Grothendieck, distincte du corps du texte.
\newcommand{\marginal}[1]{\par\smallskip\noindent\hspace{1em}%
  {\small\color{gray!60!black}#1}\par\smallskip}

% --- Titre ------------------------------------------------------------------

\AtBeginDocument{%
  \begin{center}
    {\large\bfseries Fonds Alexandre Grothendieck --- cote n\textsuperscript{o}~\grothFolder}\\[2pt]
    {\itshape\grothFoldertitle}\\[4pt]
    {\small feuillets \grothFirst--\grothLast\ (lot \grothBatch)}\\[2pt]
    {\footnotesize\color{gray!60!black}Transcription établie d'après le fac-similé numérisé
     par l'Université de Montpellier.\\
     Les lectures incertaines sont soulignées, les illisibles notées [\,\ldots\,],
     les ajouts entre crochets.}
  \end{center}
  \vspace{1em}}

\endinput


\folder{161-2}
\batch{2}
\pages{21}{40}
\dating{[vers 1963-1973]}
\watermark{Édition de démonstration}
\foldertitle{Algèbre universelle [ou catégories] : notes manuscrites (s.d.)}

\begin{document}

\page{21}
\note{la page reprend en cours de phrase la construction laissée p. 20}
des produits $2$-fibrés dans $(\mathrm{Top})$, [\struck{des} $R \times^{2}_{S} R_{1}$],
i.e. $2$-limites amalgamées dans $\Delta$ : $R \amalg_{S} R_{1}$. Cette construction
est $\pm$ connue, redonnons-en une autre, dans le cas général $R \times^{2}_{T} S$,
avec les foncteurs images inverses
\[
  T \xrightarrow{\ f^{*}\ } R, \qquad T \xrightarrow{\ g^{*}\ } S .
\]
On choisit des \struck{familles} \add{petites sous-catégories} génératrices
$T_{0}, R_{0}, S_{0}$ telles que $f^{*}(T_{0}) \subset R_{0}$,
$g^{*}(T_{0}) \subset S_{0}$. On prend $U_{0}$ la \add{catégorie} $2$-somme
amalgamée de $R_{0}$ et $S_{0}$, de sorte que pour tout $(E \struck{\ill{}}) \ldots$
\[
  \underline{\mathrm{Hom}}(U_{0}, E)
  \;\xrightarrow{\ \approx\ }\;
  \underline{\mathrm{Hom}}(R_{0}, E)
    \times^{2}_{\underline{\mathrm{Hom}}(T_{0},E)}
  \underline{\mathrm{Hom}}(S_{0}, E)
\]
Il faut, si $E \in \mathrm{Ob}\,\Delta$, exprimer qu'un objet des \struck{seconds
mb.} premiers mb., $U_{0} \xrightarrow{h} E$, est dans \struck{\ill{}} la
sous-catégorie pleine
\[
  \underline{\mathrm{Hom}}_{\Delta}(R, E)
    \times^{2}_{\underline{\mathrm{Hom}}_{\Delta}(T,E)}
  \underline{\mathrm{Hom}}_{\Delta}(S, E)
\]
des derniers mb. Si on a pris $R_{0}, S_{0}$ stables par $\varprojlim$ finies,
cela s'exprime aisément par la \struck{commutativité} \struck{de la aux}
condition que $h$ transforme certains cônes projectifs finis de $U_{0}$
(\struck{provenant} \add{aux} cônes exacts \struck{de} $R_{0}$ et de $S_{0}$)
en cônes exacts, et \struck{que les uns} qu'elle transforme certaines familles
(provenant \ill{} familles couvrantes de $R_{0}$ et $S_{0}$) en familles
couvrantes. Donc on réalise la théorie envisagée
\[
  T = T_{R} \times^{(2)}_{T_{T}} T_{S}
\]
comme sous-théorie pleine (classifiée par des axiomes d'\struck{\ill{}}
\ill{}, finales) de $T_{\rho_{\Delta}(U_{0})}$, et qui donne la conclusion.
Mais sachant désormais que \struck{dans} la théorie $T$ est représentable par
un topos, \struck{et} revenant à la description générale comme sous-théorie
d'un $T_{\rho_{\Delta}(U_{0})}$ (\add{plus la} \struck{sans} condition que
$R_{0}, S_{0}$ stables par $\varprojlim$ exactes), on voit que c'est donné par
un \underline{sous-topos} de $\rho_{\Delta}(U_{0})$. Ainsi on trouve, si
$T_{0} \to S_{0}$ est surjectif sur les objets (comme dans les cas qui nous
intéressaient) que $R_{0} \to U_{0}$ est surjectif sur les objets, donc
$R \to U = R'$ \struck{est} \struck{généré} \ill{} une image qui est génératrice.

\page{22}
\textbf{d) \underline{Passage à la limite}}\quad Cela \ill{} comme suit : les
$2$-limites \add{inductives} \struck{stables} dans $\Delta$, il \struck{faut}
\ill{} que les $2$-limites transfinies existent et soient des topos. La
démonstration donnée dans ce monde en fait, par le système inductif
\struck{de topos} dans $\Delta$ formé de topos [des systèmes projectifs dans
$(\mathrm{Top})$] $\varepsilon(\bar{R}_{\alpha})$ : on se débrouille pour
choisir des \struck{systèmes de} \add{petites cat.} génératrices
\add{\uncertain{Cocha}} \struck{qui} soient stables par les morphismes de
transition \struck{\ill{}} — p. ex. à l'aide des filtrations cardinales, mais
c'est un mauvais \ill{} un monde — et on prend
\struck{la catégorie $2$-\ill{}} la
\[
  C = \varinjlim_{\alpha,\,(\mathrm{Cat})} C_{\alpha},
\]
on peut donner un sous-topos de $\hat{C}$.

En l'occurrence, \add{dans} \struck{le cas filtrant} \struck{c'est que}
on veut que si \struck{un} $R$ est le topos limite, \struck{et si} les images
dans $R$ des $R_{\alpha}$ forment un syst. de générateurs
\struck{\ill{} générateurs} \add{d'où} c'est une syst. \struck{génératrice} :
et si $J$ s'envoyant en le syst. inductif des $R_{\alpha}$, $J \to R_{\alpha}$
\struck{engendre} $R_{\alpha}$ pour tout $\alpha$, alors $J \to R$ engendre $R$.

\note{annotations portées au-dessus du diagramme : $\rho'(I) = \hat{J_{0}}$ ;
$\rho'(J_{1})$, $\rho'$ l'\ill{}adjoint, i.e. $\Delta_{R}$ cat. avec produits
finis ; $= \rho'(J_{2})$ au-dessus de $R_{2}$}

\begin{tikzcd}[column sep=small, row sep=small, nodes={font=\scriptsize}]
R_{0} \arrow[r] & R_{1} \arrow[r, "\mathrm{loc}"] & R_{2} \arrow[r, "\mathrm{loc}"] & R_{3} = R \\
J_{0} \arrow[u] \arrow[r, "\text{bij. sur objets}"] & J_{1} \arrow[u] \arrow[r, "\mathrm{loc.}"] & J_{2} \arrow[u] \arrow[ur] & \\
I \arrow[u] & & &
\end{tikzcd}

\note{$J_{1} \to J_{2}$ porte en outre « (bij. sur ob.) » ; sous $J_{1}$ :
« $=$ analyseur engendré par $J_{0}$ (analyseur \struck{libre}) et les flèches
données \ill{} analyseur \ill{} » ; sous $J_{2}$ : « $=$
\uncertain{sous-analyseur} plein de $R$ engendré par $I \to R$, ens.
sous-jacent $J_{0}$ »}

\begin{tikzcd}[column sep=small, row sep=small, nodes={font=\scriptsize}]
\hat{J_{1}} \arrow[r, "\mathrm{loc.}"] \arrow[d, "\mathrm{loc}"'] & \hat{J_{2}} \arrow[r] \arrow[d, "\mathrm{loc}"] & R \\
R_{1} \arrow[r, "\mathrm{loc}"] & R_{2} \arrow[r, "\mathrm{loc}"] & R_{3}
\end{tikzcd}

\marginal{$R$ apparaît comme une catégorie des fractions de \struck{\ill{}}
$\hat{J_{1}}$, i.e. $J_{1}$ est l'analyseur libre sur $I$ engendré par les
flèches induites (sur $J_{0}$) par $R$, $\hat{J_{1}}$}

\page{23}
\textbf{\underline{Structure représentée par} $B_{G} = \hat{C}$}\quad
$C = \struck{\ill{}}(\cdot\,, G)$. \emph{Relativisation dans un topos.}

\textbf{\underline{Topos produit.}}\quad Remarquons que les
\[
  T(E) \simeq \underline{\mathrm{Hom}}\,\mathrm{top}(E, R)^{\circ}
        \simeq \underline{\mathrm{Hom}}_{\Delta}(R, E)
\]
\ldots ces catégories sont stables par $\varinjlim$ \underline{filtrantes}
(car il vient à peu près des conditions sur des structures : elles \ill{}
convergent \ill{} pas à des théories algébriques sur \struck{\ill{}}
\add{quelque} soit $\Delta' \supset \Delta$).

\textbf{N.B.}\quad En général on aura un autre type de lim.
\underline{Question} : $\forall E$, $T(E)$ admet des \struck{produits fibrés}
\add{$\varprojlim$ finies} quelconques (et les $T_{R}(E) \to T_{R}(E')$ à y
commuter) — auquel cas $T$ = restriction à $\Delta$ d'une théorie algébrique
sur $\Delta_{0}$, i.e. $R \simeq \hat{C}$, avec $C \subset (E \in
\mathrm{Ob}\,\Delta_{0}$ \add{(vérité)}, i.e. $C$ stable par $\varprojlim$
finies ? Si $T$ est un type, on sait que $T(E)$ est stable par $\varprojlim$
finies quelconques et accessible, \struck{\ill{}}, \add{stable} par \ill{}
générale.
\[
  T_{R}(\mathrm{Ens}) \simeq \underline{\mathrm{Pt}}(R)^{\circ}
    \simeq \underline{\mathrm{Fib}}(R) .
\]
\note{le glyphe lu $\mathrm{Pt}$ est peu net ; le contexte — points d'un topos,
et la comparaison avec $\underline{\mathrm{Fib}}$ — le soutient} \ill{} et
$R \neq \mathrm{Top}(\varphi)$.

\textbf{\underline{Sous-topos du topos disons \add{$R$} engendré par un objet
$X$}} (a priori de structure \ill{} sur un objet de \ill{}) : comparer avec un
$f$-\ill{}, avec $R = \hat{C}$, $C$ = catégorie opposée \add{[pts à les
$\Delta_{n}$]} aux ens. finis, donc $R$ = ens. semi-simpliciaux. Les
\underline{sous-topos} \add{(ouverts)} \ill{} de $C$ ; on trouve les
sous-topos suivants (classés \ill{} par leurs points) :

\begin{itemize}
  \item $\{\, R,\ \emptyset,\ R^{+} = \varprojlim (X \to e) \,\}$ ;
  \item de p. f. : $\{\, \mathbb{N},\ \emptyset,\ \mathbb{N}^{+} \,\}$ ;
  \item \uncertain{(irréd.)} : $\{\, S,\ \emptyset,\ (\mathrm{Ens}\ \ill{}
    \ \ill{})^{*} \,\}$, où $S = (\mathrm{Ens}_{f})^{\circ}$ ;
  \item sous-thèmes : $\{\, \struck{T(E) = E},\ \emptyset,\ T(E) \subset E \,\}$,
    \add{sous-catégorie pleine des objets connexes de $\bar{E}$}, d'où
    $T(E) = \emptyset$ sauf si $E = \mathrm{Top}(\varphi)$.
\end{itemize}

\underline{Sous-topos fermés complémentaires} : p. h. f.
$\{\, \emptyset,\ R,\ \Phi_{0} = \mathcal{C}U \,\}$ ;
$\{\, \emptyset,\ \mathbb{N},\ \{\Delta_{p}\} \,\}$ ;
$T(E) \subset E$ \add{cat. \ill{} réduite à \struck{\ill{}} \ill{}}.

\page{24}
Il y a enfin un sous-topos fermé, \struck{maximal loc.} dont les \ill{} et
dernier — le topos $V\struck{\ill{}}$, qui n'est ni ouvert, ni \ill{}, dont les
pts \ill{} $\Delta_{0}, \Delta_{1}$, et qui \struck{\ill{}} :
\[
  T_{V}(E) = \text{End. des sous-objets de } \struck{\ill{}}\, e_{E}
\]
(i.e. un objet $X \in \mathrm{Ob}\,E$ tel que $x \to x_{e}$ mono.), et le
deuxième topos $\mathrm{Ens}^{+}_{\wedge} V$ \struck{\ill{}} : le plus \ill{},
\add{plus} \ill{} $\Delta_{1}$, tel que
\[
  T_{\mathrm{Unv}}(E) \subset E, \qquad \text{i.e. } x \xrightarrow{\ \sim\ } e_{E}
\]
\add{objets finaux}.

En résumé \ill{} : $6$ sous-topos, avec diagramme d'implications

\begin{tikzcd}[column sep=small, row sep=small, nodes={font=\scriptsize}]
 & R & & \\
R^{+} \arrow[ur, no head] & & V = \overline{(\varphi_{0},\varphi_{1})} \arrow[ul, no head] & \\
 & R^{+} \wedge V = \varphi_{1} \arrow[ul, no head] \arrow[ur, no head] & & \overline{(\varphi_{0})} \arrow[ul, no head] \\
 & & \emptyset \arrow[ul, no head] \arrow[ur, no head] &
\end{tikzcd}

\note{la moitié droite du feuillet porte une première version du même treillis,
biffée en entier ; elle n'est pas transcrite}

\textbf{N.B.}\quad \struck{\ill{}} $V$ et $R^{+} \wedge V$ ne sont pas des
complémentaires, mais ils ont les complémentaires faibles $\emptyset$ et
$\varphi_{0}$. $V$ \struck{représent\ill{}} \ill{} représente le foncteur
\[
  E \longmapsto \mathcal{P}(E) = \mathcal{P}(e_{E}) .
\]
\note{le glyphe est un $P$ barré ; rendu ici $\mathcal{P}$, l'objet des parties}

\textbf{N.B.}\quad

\begin{itemize}
  \item Pas d'objets loc. constants, ou constants \add{pendants} \ill{} $\varphi$ ;
  \item pas d'objets \struck{constants} \ill{} \add{quelconque} sur principe
    cela \struck{\ill{}} (resp. si sur $e$) ;
  \item pas d'objets constructibles.
\end{itemize}

\textbf{\underline{Question}}\quad Topos classifiant et sous-topos pour les
structures : $f : X \to Y$ sans plus, où $f : X \to X$ sans plus
[$= B_{\mathbb{N}^{\times}}$], $f : e \to X$,
\[
  X \times Y \to Z, \qquad X \times Y \to \struck{X},
\]
\[
  X \times X \longrightarrow X \quad (!), \qquad
  X \times X \rightrightarrows X \quad (!!) .
\]

\marginal{\ill{} — note portée verticalement dans la marge gauche, en regard de
« $X, Y$ »}

\page{26}
\note{le théorème est repris ici depuis le début, sous une forme plus générale
que celle de la p. 3 : les cônes y sont indexés par un ensemble $\sigma$ de
couples et un ensemble $\rho$ de systèmes}

\textbf{\underline{Th.}}\quad Soit

\begin{enumerate}
  \item[1)] $\sigma$ un ensemble de couples $(S, \lambda)$, $S$ catégorie,
    $\lambda = (\lambda', \lambda'')$ ens. de cônes (proj. et ind.) dans $S$ ;
  \item[2)] $\rho$ un ens. de \add{systèmes} \struck{diagrammes}
    $(\struck{\ill{}}, R_{1}, \delta_{1}, R_{2}, \delta_{2},
    \varphi\,\struck{\ill{}})$, où les $\delta_{i} = (\delta'_{i},
    \delta''_{i})$ $(i = 1, 2)$ sont des ens. de cônes, $R_{i}$ une catégorie,
    \struck{une catégorie avec des \ill{} foncteurs \ill{}} et un foncteur
    $R_{2} \xrightarrow{\varphi} R_{1}$ tel \struck{$R_{1} \to R_{2}$} que
    $\varphi(\delta_{2}) \subset \delta_{1}$
    \struck{$(R_{1}) \longrightarrow (\delta_{2})$}.
\end{enumerate}

Soit $\Delta = \Delta_{\sigma,\rho}$ la \add{sous-}catégorie de
$\underline{\mathrm{Cat}}$ définie ainsi :

\begin{enumerate}
  \item[a)] les objets sont les catégories $E$ ayant les propriétés suivantes :
  \item[a$_{1}$)] $\forall (S,\lambda) \in \sigma$ et tout foncteur
      $S \xrightarrow{u} E$, les systèmes projectifs et inductifs
      \ill{} aux cônes $\in u(\lambda)$ ont une limite ;
  \item[a$_{2}$)] $\forall (\struck{\ill{}}, R_{1}, \delta_{1}, \ldots) \in
      \rho$, \struck{et tout foncteur $u_{1} : R_{1} \to E$ qui \ill{}
      $\delta_{1}$-ex. ct, il existe un foncteur $u_{2} : R_{2} \to E$ faisant
      $\delta_{2}$-exact et tel que $u_{1}\varphi_{1} \to u_{2}\varphi_{2}$},
      \add{$u \mapsto u \circ \varphi$ induisant une équivalence} :
      \[
        \underline{\mathrm{Hom}}_{\delta_{1}}(R_{1}, E)
        \;\xrightarrow{\ \approx\ }\;
        \underline{\mathrm{Hom}}_{\delta_{2}}(R_{2}, E) ;
      \]
  \item[b)] pour $E, E' \in \mathrm{Ob}\,\Delta$,
    $\underline{\mathrm{Hom}}_{\Delta}(E, E')$ est la sous-catégorie pleine
    \add{(stab.)} de $\underline{\mathrm{Hom}}_{\mathrm{Cat}}(E, F)$ formée des
    foncteurs $F : E \to E'$ tels que l'on ait : $\forall (S, \lambda) \in
    \sigma$ et tout foncteur $u : S \to E$ qui est $\lambda$-ex. ct, $F \circ u$
    est $\lambda$-exact. (Il est clair que cette condition est stable par
    isomorphisme et par composition, et vérifiée \struck{par les} \add{si $F$
    est une} équivalence de catégories.)
\end{enumerate}

On suppose satisfaites les conditions

\begin{enumerate}
  \item[A)] $\sigma$ et $\rho$ sont petits, \struck{les} $\forall (S,\lambda)
    \in \sigma$ $S$ soit petit et $\lambda$ soit petit, et les
    \add{catégories} \struck{diagrammes} d'indices des cônes $\in \lambda$
    soient petites ; de même pour dans $(\struck{\ill{}}, R_{1}, \delta_{1},
    R_{2}, \delta_{2}, \varphi\,\struck{\ill{}}) \in \rho$ les catégories
    $\struck{\ill{}}\,R_{i}$ $(i = 1, 2)$ soient petites, les $\delta_{i}$
    soient petits et les catégories d'indices des cônes $\in \delta_{i}$
    soient petites.
  \item[B)] $\forall (R, \ldots) \in \rho$, $(R_{i}, \delta_{i}) \in \sigma$
    pour $i = 1, 2$. [On le veut \add{la condition d'exactitude de cet et} (ou
    du moins : \struck{\ill{}} le foncteur défini par $(R_{i}, \delta_{i})$
    \struck{est conséquence des} \add{celles définies par $\sigma$}).]
    Introduire des conditions d'exactitude supplémentaires b).
\end{enumerate}

\page{27}
\ldots sur les catégories envisagées pour $\Delta$ (au plus des conditions
d'existence \add{$E$} au lieu dans a), que si les \add{types de}
(\underline{lim} que les formulations des conditions impliquent existant
\underline{tjs} dans les $E$ envisagées, et si les foncteurs admis entre ces
$E$ y commutent\ldots)].

Soit de plus $C$ une catégorie ess. petite, $\delta$ \struck{diag.}
\add{cônes} $= (\delta', \delta'')$ un petit ensemble de \ill{} (proj. et ind.)
\struck{définissant} correspondant à de petites catégories d'indices.

Alors il existe un $\underline{C' \in \mathrm{Ob}\,\Delta}$ et un foncteur
$i : C \to C'$ qui est $\underline{\delta\text{-exact}}$, de telle façon que
$(C', i)$ soit $\underline{2\text{-universel}}$, i.e. pour tout $E \in
\mathrm{Ob}\,\Delta$, le foncteur $F \mapsto F \circ i$ induit une équivalence
\[
  \underline{\mathrm{Hom}}_{\Delta,\, i(\delta)}(C', E)
  \;\xrightarrow{\ \approx\ }\;
  \underline{\mathrm{Hom}}_{\mathrm{Cat},\, \delta}(C, E) .
\]
De plus, cette catégorie $C'$ (définie à équivalence \struck{\ill{}} près
d'après i) est ess. petite.

\textbf{\underline{Dém.}}\quad On construit par réc. transfinie un
\add{(ps) syst. inductif} \struck{suite} \add{chaîne} \ill{} finie
$(C_{i}, \delta_{i})$ de catégories sur $C$, munies d'un $\delta_{i}
\add{= (\delta'_{i}, \delta''_{i})}$ \struck{de} cônes, en procédant ainsi :
\begin{gather*}
  (C_{0}, \delta_{0}) = (C, \delta) \\
  (C_{i+1}, \delta_{i+1}) = K(C_{i}, \delta_{i}), \qquad
  \tau_{i} : C_{i} \to C_{i+1} \\
  \text{i.e. } K \text{ est une opération à expliciter plus bas} \\
  C_{i} = \varinjlim_{j < i,\,(\mathrm{Cat})} C_{j}, \qquad
  \delta_{i} = \bigcup_{j < i}
    \bigl(\mathrm{Im}\ \text{de } \delta_{j} \text{ par } C_{j} \to C_{i}\bigr)
  \quad \text{si } i \text{ ordinal limite.}
\end{gather*}
On note $\alpha_{i} : C \longrightarrow C_{i}$.

La construction $K$ (indépendante de $i$) a les propriétés suivantes, que pour
la commodité des références nous écrivons avec les notations
$(C_{i}, \delta_{i})$, $(C_{i+1}, \delta_{i+1})$ :

\page{28}
\begin{enumerate}
  \item[a)] $\forall E \in \mathrm{Ob}\,\Delta$, le foncteur $u \mapsto u \circ
    \tau_{i}$ induit une équivalence
    \[
      \underline{\mathrm{Hom}}_{\delta_{i+1}}(C_{i+1}, E)
      \longrightarrow
      \underline{\mathrm{Hom}}_{\delta_{i}}(C_{i}, E) ;
    \]
  \item[b)] $\displaystyle \tau_{i}(\delta_{i}) \subset \delta_{i+1} \subset
    \tau_{i}(\delta_{i}) \cup \bigcup_{(S,\lambda) \in \sigma} \;
    \bigcup_{u \in \underline{\mathrm{Hom}}(S, C_{i+1})} u(\lambda)$ ;
  \item[c)] $\tau_{i} : C_{i} \to C_{i+1}$ est \underline{préexact} par
    \struck{les} $\delta_{i}$, i.e.
\end{enumerate}

\note{suit un alinéa entièrement biffé de traits diagonaux, encore lisible :
« $\forall (\struck{\ill{}}, R_{1}, \delta_{1}, R_{2}, \delta_{2},
\varphi\,\struck{\ill{}}) \in \rho$, et $u_{1} : R_{1} \to C_{i}$
\add{tel que $u_{1}(\delta_{1}) \subset \delta_{i}$}, il existe un foncteur
\struck{$\delta$-exact} $v_{2} : R_{2} \to C_{i+1}$ tel que $v_{2} \circ
\varphi_{2} = (\tau_{i} u_{1}) \circ \varphi_{1}$ \add{(et $v_{2}(\delta_{2})
\subset \delta_{i+1}$)} » ; il est repris en e$_{0}$)--e$_{2}$)}

\begin{enumerate}
  \item[d)] $\forall (S, \lambda) \in \sigma$ \struck{et} $u : S \to C_{i}$,
    \struck{et une \struck{diag.} catégorie $I'$ ($I''$)}
    \struck{intérieurement dans $\in \lambda'$ ou $(\lambda'')$}
    \add{(i.e. $\lambda = (\lambda', \lambda'')$)}, il existe
    \struck{un cône $c'$ sur le syst. proj. \ill{} $(i-1)$ correspondant à
    $\tau_{i}(c)$} \add{$c' \in \delta_{i+1}$} un foncteur
    $\bar{v} : \bar{S} \to C_{i+1}$ (où $\beta : S \to \bar{S}$ est obtenu en
    dédoublant les \ill{} \add{sommets} des cônes critiques de $S$ (**)) tel
    que $\bar{v}\beta = \tau_{i} u$, et tel que $\bar{v}(\bar{\lambda}) \subset
    \delta_{i+1}$ ;
  \item[e$_{0}$)] $\forall (R_{1}, \delta_{1}, \ldots) \in \rho$ et
      $u_{1}, v_{1} : R_{1} \rightrightarrows C_{i}$ \add{tels que
      $u_{1}(\delta_{1}) \subset \delta_{i}$, $v_{1}(\delta_{1}) \subset
      \delta_{i}$}, et $\alpha, \beta : u_{1} \rightrightarrows v_{1}$, la
      relation $\alpha * \varphi = \beta * \varphi$ implique
      $\tau_{i} * \alpha = \tau_{i} * \beta$ ;
  \item[e$_{1}$)] $\forall (R_{1}, \delta_{1}, \ldots) \in \rho$ et
      $u_{1}, v_{1} : R_{1} \rightrightarrows C_{i}$ tels que
      $u_{1}(\delta_{1}), v_{1}(\delta_{1}) \subset \delta_{i}$, et tout hom.
      $u_{1} \circ \varphi \xrightarrow{\ \alpha_{2}\ } v_{1} \circ \varphi$,
      $\exists$ un hom. $\tau_{i} u_{1} \xrightarrow{\ \bar{\alpha}_{1}\ }
      \tau_{i} v_{1}$ tel que $\bar{\alpha}_{1} * \varphi = \tau_{i} *
      \alpha_{2}$ ;
  \item[e$_{2}$)] $\forall (R_{1}, \delta_{1}, \ldots) \in \rho$ et
      $u_{2} : R_{2} \to C_{i}$ tel que $u_{2}(\delta_{2}) \subset \delta_{i}$,
      $\exists\ \bar{u}_{1} : R_{1} \to C_{i+1}$ tel que $\bar{u}_{1} \circ
      \varphi = \tau_{i} \circ u_{2}$ et $\bar{u}_{1}(\delta_{1}) \subset
      \delta_{i+1}$.
\end{enumerate}

\marginal{(*) prouver que pour $\omega$ grand, $(C_{\omega}, \tau_{\omega})$
\ill{} ; et cf. \ill{} (**) plus bas}
\note{les renvois (*) et (**) sont développés p. 31}

\page{29}
De a) on déduit par récurrence transfinie, pour $\forall E \in
\mathrm{Ob}\,\Delta$
\[
  (1)\qquad
  \underline{\mathrm{Hom}}_{\delta_{i}}(C_{i}, E)
  \;\xrightarrow{\ \approx\ }\;
  \underline{\mathrm{Hom}}_{\delta}(C, E)
  \qquad (F \mapsto F \circ \alpha_{i}) .
\]
De b) on déduit de même
\[
  (2)\qquad
  \alpha_{i}(\delta) \subset \delta_{i} \subset
  \alpha_{i}(\delta) \cup
  \bigcup_{\substack{j < i \\ (S,\lambda) \in \sigma}} \;
  \bigcup_{u \in \underline{\mathrm{Hom}}(S, C_{j})} \tau_{ij}\, u(\lambda)
  \qquad \text{si } i \text{ ordinal limite.}
\]

Prouvons que pour $\omega$ \add{ordinal limite} assez grand ($\omega$ un
« petit » ordinal), $(C_{\omega}, \delta_{\omega})$ est une solution du Pb.
Cela \struck{signifie} \add{Il suffit} qu'on ait

\begin{enumerate}
  \item[1\textsuperscript{o})] $C_{\omega} \in \mathrm{Ob}\,\Delta$, i.e.
  \item[A$_{1}$)] $\forall (S,\lambda) \in \sigma$ et $u : S \to C_{\omega}$,
      existence des lim. correspondant aux cônes $\in u(\lambda)$ ;
  \item[A$_{2}$)] $\forall (\struck{\ill{}}, R_{1}, \ldots, \varphi) \in
      \rho$ \struck{et \ill{}},
      $\underline{\mathrm{Hom}}_{\delta_{1}}(R_{1}, C_{\omega})
      \xrightarrow{\ \sim\ }
      \underline{\mathrm{Hom}}_{\delta_{2}}(R_{2}, C_{\omega})$ ;
  \item[2\textsuperscript{o})] $\alpha_{\omega} : C \to C_{\omega}$ est
    $\delta$-exact ;
  \item[3\textsuperscript{o})] $\displaystyle \struck{\ill{}}\ \ \delta_{\omega}
    = \alpha_{\omega}(\delta) \cup \bigcup_{(S,\lambda) \in \sigma} \;
    \bigcup_{u \in \underline{\mathrm{Hom}}_{\lambda}(S, C_{\omega})}
    u(\lambda)$.
\end{enumerate}

[Car 3\textsuperscript{o} implique que pour $\forall E \in \mathrm{Ob}\,\Delta$
on a
\[
  \underline{\mathrm{Hom}}_{\Delta,\, \alpha_{\omega}(\delta)}(C_{\omega}, E)
  \ \struck{\ill{}} = 
  \underline{\mathrm{Hom}}_{\struck{\ill{}},\, \delta_{\omega}}(C_{\omega}, E)
  \;\xrightarrow[\ (1)\ ]{\ \approx\ }\;
  \underline{\mathrm{Hom}}_{\delta}(C, E) .]
\]

\struck{Alors} Soit $I_{\omega}$ l'ensemble des ordinaux $j < \omega$. Pour
vérifier 2\textsuperscript{o}) il suffit de vérifier l'assertion plus forte
\[
  (3)\qquad \forall i < \omega, \quad
  \tau_{\omega,i} : C_{i} \to C_{\omega} \text{ est } \delta_{i}\text{-exact.}
\]
Ceci résulte à son tour de c) si $\omega$ est grand devant les cardinaux : soit
$F \in I$, où $I$ est une cat. d'indices intervenant \add{dans} les cônes $\in
\delta_{i}$ ; or en vertu de (2), \struck{les $I$ sont \ill{}} \add{de plus}
ceux qui interviennent dans $\delta$, soit dans un $\lambda$ associé à un
$(S, \lambda) \in \sigma$ — c'est une \ill{} \ill{} des petites catégories,
donc les $\mathrm{ord}\, F \in I$ sont majorés par un cardinal $\Pi$ — et il
suffit \struck{$\tau_{\omega}$} \add{$\omega$} grand devant $\Pi$.

\page{30}
Il résulte de (3) et (2) appliqué à $i = \omega$ que l'on a
\[
  \alpha_{\omega}(\delta) \subset \delta_{\omega} \subset
  \alpha_{\omega}(\delta) \cup \bigcup_{(S,\lambda) \in \sigma} \;
  \bigcup_{u \in \underline{\mathrm{Hom}}_{\lambda}(S, C_{\omega})}
  \struck{\ill{}}
\]
\marginal{mod. isom. des cônes}
et pour prouver 3\textsuperscript{o}) il faut donc prouver

\begin{enumerate}
  \item[4)] $\forall (S,\lambda) \in \sigma$ et $u \in
    \underline{\mathrm{Hom}}_{\lambda}(S, C_{\omega})$, \struck{exi}$\exists 
    i < \omega$ et un hom. $u_{i} : S \to C_{i}$ tel que $u_{i}(\lambda)
    \subset \delta_{i}$ (et $u \simeq \tau_{\omega,i} \circ u_{i}$
    \add{(mod. is. des cônes)}).
\end{enumerate}

\note{suivent deux lignes biffées reprenant l'argument par 1\textsuperscript{o})}

\begin{enumerate}
  \item[a$_{1}$)] sera conséquence de \struck{d)} si nous notons que tout
    foncteur $u : S \to C_{i}$ se factorise par un $C_{i}$ — pourvu qu'on
    prenne $\Pi \geq \mathrm{card}\, F \in S$ pour tout $(S,\lambda) \in
    \sigma$ ;
  \item[a$_{2}$)] sera conséquence de e), utilisant B (les $(R_{i},
    \delta_{i}) \in \sigma$) et $\beta$ \struck{(4)} \add{(3)} au lieu de i).
\end{enumerate}

\note{une ligne entière biffée : « Donc il reste à donner \ill{} des conditions
dans la construction $\alpha$ qui assurent (3), et à expliciter une
construction $K$ » ; en marge, de sa main, « peut remonter ! »}

Prouvons 4). On sait que $u : S \to C_{\omega}$ se factorise \struck{par}
$u_{i} : S \to C_{i}$ par un $C_{i}$ (pourvu qu'on prenne $\Pi \geq
\mathrm{card}\, F \in S$ $\forall (S,\lambda) \in \sigma$), donc quitte à
composer avec $\tau_{i}$, on peut supposer que ceci se fait \ill{} par
$\bar{u}_{i} : \bar{S} \to C_{i}$ tel que $\bar{u}(\bar{\lambda}) \subset
\delta_{i}$, donc $\tau_{\omega,i} \circ \bar{u}_{i}$ est un hom.
$\bar{S} \to C$ qui est $\bar{\lambda}$-exact, \add{qu'on a (3)}.
\struck{\ill{}} Comme $\bar{u}(c_{c})$ \ill{} $\bar{u}_{i}(i_{c})$ donnent un
isom. pour $\forall c \in \lambda$, donc \ill{} $\bar{u}_{j}(i_{c})$ \ill{}
pour $\forall c \in \lambda$, donc les $\tau_{ji}\,\bar{u}_{i}(i_{c})$
\ill{} propriété pour $j \in I_{\omega}$ assez grand (pourvu que
$\Pi \geq \mathrm{card}\,\lambda = \mathrm{card}\{c\}$), donc
\struck{\ill{}} les \struck{éléments} cônes des $u_{j}(\lambda)$ sont isom. à
ceux des $\bar{u}_{j}(\bar{\lambda})$ ; comme ceux-ci sont dans $\delta_{j}$,
il en est de \ill{} de ceux-ci, à isom. près, cqfd.

\page{31}
\note{la page porte les deux renvois annoncés p. 28}

(**) Soit $S$ une catégorie avec $\lambda = (\lambda', \lambda'')$ ens. de
cônes. On lui associe une catégorie $\bar{S}$ et un foncteur
$S \xrightarrow{\ \beta\ } \bar{S}$ \struck{tel que} en ajoutant, pour
$\forall$ \struck{sommet de} cône
$c' : \bigl(x \xrightarrow{\ p_{\alpha}\ } (x_{\alpha})\bigr) \in \lambda'$
\struck{\ill{}}, un objet $\bar{x}_{c}$ à $S$, avec \struck{une de} une
famille de flèches $x \to \bar{x}_{c}$ et \add{une famille de flèches}
$\bar{x}_{c} \xrightarrow{\ q_{\alpha}\ } x_{\alpha}$ factorisant les
$p_{\alpha}$ — et de manière duale pour les cônes $c'' \in \lambda''$.

\begin{tikzcd}[column sep=small, row sep=small, nodes={font=\scriptsize}]
 & x_{1} & \\
x \arrow[r, "i_{c}"] & \bar{x}_{c} \arrow[u, "q_{1}"] \arrow[r, "q_{2}"] \arrow[d, "q_{3}"'] & x_{2} \\
 & x_{3} &
\end{tikzcd}

\note{le croquis marginal porte en outre les flèches $p_{\alpha} : x \to
x_{\alpha}$, qui se factorisent par $i_{c}$}

(On dit que $(\bar{S}, \bar{\lambda})$ est \underline{déduit} de $(S, \lambda)$
en \underline{dédoublant} les sommets des cônes critiques de $S$ ;
$\bar{\lambda}$ est formé des cônes déduits des cônes $c \in \lambda$ en
remplaçant le sommet $x_{c}$ par $\bar{x}_{c}$.) Donc pour $\forall$ foncteur
$u : S \to E$, avec $E \in \mathrm{Ob}\,\Delta$, il existe une
\underline{unique} extension en $\bar{u} : \bar{S} \to E$ qui soit
$\bar{\lambda}$-exacte ; de plus, $u$ est $\lambda$-ex. $\Leftrightarrow$
\struck{ses \ill{}} pour $\forall c \in \lambda$, $\bar{u}(i_{c})$ est un
isomorphisme.

(*) Définition du foncteur $\underline{\delta\text{-préexact}}$
$D \xrightarrow{\ u\ } F$, $D$ muni d'un $\delta = (\delta', \delta'')$ de
cônes\ldots

\page{33}
\note{feuillet de travail rapide, en anglais et en français mêlés, sans ordre
linéaire : les fragments sont donnés ici dans l'ordre où ils tombent sur la
page, du haut vers le bas}

\[
  H^{n}(K(G,1), \Pi) \xrightarrow{\ \sim\ } \mathrm{Simpl}(\mathrm{ENS}),
  \qquad \struck{Z_{i}(\Pi)} .
\]

\emph{Lawvere theory}\quad $\mathcal{L}$ a small cat., $\mathbb{1} \in
\mathcal{L}$, $n = \mathbb{1}\ldots\mathbb{1}$.
\[
  S = \underline{\mathrm{Hom}}'(\mathcal{L}^{\circ}, \mathrm{ENS})
  \ \rightleftarrows\ \mathrm{Alg}(\mathfrak{T}),
  \qquad (\mathcal{L}^{\circ} \to \mathrm{ENS}), \qquad
  \mathrm{Mod}(\mathcal{L}^{\circ}) .
\]
avec $F$ adjoint à gauche de $U$, et $(\mathrm{ENS}) \longrightarrow \ill{}$.

\[
  G_{1} \rightrightarrows B_{2} \dashrightarrow S * S_{1} .
\]

$P \in \mathrm{Ob}(S)$ :

\begin{enumerate}
  \item[1\textsuperscript{o}] $P_{I} = \coprod_{I} P$, \quad $P \times \Pi$ ;
  \item[2\textsuperscript{o}] \emph{Every eq. rel is effective},
    $R \rightrightarrows S$ ;
  \item[3\textsuperscript{o}] $f : X \to Y$ \underline{effective}
    $\Longleftrightarrow$ $\mathrm{Hom}(P, X) \to \mathrm{Hom}(P, Y)$
    surjectif \add{(epi)} ;
  \item[4.] \underline{fiber prod} ;
  \item[5\textsuperscript{o}] $P \to P \times K$, $P \times n$,
    $\mathrm{Hom}(P, \Pi)$.
\end{enumerate}

\emph{morphism of \struck{the} \underline{effective} descent.}

\note{le fragment central est la reconnaissance d'un foncteur tripleable}
\[
  A \xrightarrow[\ \sim\ ]{\ \Phi\ } \mathrm{Alg}(\mathbb{T}),
  \qquad F \dashv U : A \to B, \qquad B \xrightarrow{\ T\ } B,
  \qquad \mathbb{T} = (UF, -, -) .
\]
\uncertain{Beck}, \uncertain{Barr}, \underline{Linton}. \quad $U$
\emph{tripleable}.
\note{les trois noms sont peu nets ; le sujet du feuillet — la tripleabilité de
$U$ — les soutient}

\begin{enumerate}
  \item[1\textsuperscript{o}] $U$ \emph{is} \struck{conservative}
    \add{conserve} ;
  \item[2\textsuperscript{o}] \ill{}
\end{enumerate}

\[
  X_{1} \rightrightarrows X_{0} \longrightarrow Q' \ \text{ in } A,
  \qquad
  U(X_{1}) \rightrightarrows U(X_{0}) \rightleftarrows Q, \qquad U(Q') ,
\]
\[
  \Longrightarrow \quad X' \rightrightarrows X \rightrightarrows Y .
\]

\[
  T^{2} \ \substack{\xrightarrow{\ T(\xi)\ } \\ \xrightarrow{\ \mu\ }}\ TX
  \xrightarrow{\ \xi\ } X, \qquad \eta, \qquad T(X) \xrightarrow{\ \xi\ } X .
\]
\note{les identités simpliciales notées à gauche : $p_{0}s_{1} = s_{0}p_{0}$ ;
$p_{1}s_{1} = \mathrm{id}$ ; $p_{0}s_{0} = \mathrm{id}$ ; en regard le diagramme
$X_{1} \rightrightarrows X_{0} \rightarrow X_{-1}$}

\page{34}
\note{feuillet de travail comme le précédent ; le sujet est le critère de
descente pour un morphisme de topos et la comonade $\pi = f^{*}f_{*}$}

\[
  X_{1} \to X, \qquad U X_{1} \to U X, \qquad C/S, \qquad C/S' .
\]
\[
  X_{1} \ \substack{\xrightarrow{\ d_{1}\ } \\ \xrightarrow{\ d_{0}\ }}\ X_{0}
  \longleftrightarrow u, \qquad X \longrightarrow X' .
\]
\[
  X_{1} \ \substack{\xrightarrow{\ d_{0}\ } \\ \xrightarrow{\ d_{1}\ }}\ X_{1}
  \xrightarrow{\ p\ } Q, \qquad Q \rightrightarrows Q ,
\]
\[
  d_{1} s_{1} = \mathrm{id}_{X_{1}}, \qquad d_{0} s_{1} = s_{0} p .
\]

$(\mathrm{Ens})^{I}$, $\mathfrak{A}$ ; $I$ \ill{} donné ;
$B = \uncertain{\mathrm{Top}}_{f}(\mathfrak{A})$.

\begin{tikzcd}[column sep=small, row sep=small, nodes={font=\scriptsize}]
 & B \arrow[dl] \arrow[dd, "f"] \\
B_{\pi} \arrow[dr] & \\
 & B'
\end{tikzcd}

\begin{tikzcd}[column sep=small, row sep=small, nodes={font=\scriptsize}]
B \arrow[d, "f_{*}", bend left=25] \\
B' \arrow[u, "f^{*}", bend left=25]
\end{tikzcd}

\[
  \pi = f^{*} f_{*} : B \longrightarrow B
  \qquad \text{(exact à g., accessible)}, \qquad
  \begin{cases} \pi \to \mathrm{id} \\ \pi \to \pi^{2} \end{cases}
\]

$f$ conservatif $\Longleftarrow$ (déf) $f^{*}$ fidèle ($=$ conservatif) — les
deux ; puis
\[
  \begin{cases}
    f^{*} \text{ conservatif} \\
    f^{*} \text{ commute aux \ill{} noyaux de couples}
  \end{cases}
\]
\[
  \mathrm{CoAlg}(B, \pi) \rightrightarrows B, \qquad
  \ill{} = B_{\pi}, \qquad
  \bigl(B \to B_{\pi}\bigr), \qquad
  B_{\pi} \xrightarrow{\ \simeq\ } B \quad \text{(th)} .
\]
\note{à droite, en regard : « $B \to B_{\pi}$ conservatif », le mot
\struck{fidèle} biffé au profit de « conservatif »}

\[
  \underline{\mathrm{Hom}}_{\mathrm{Ens}}(f^{*}X, Y)
  \simeq \underline{\mathrm{Hom}}_{G}(X, f_{*}(Y)),
  \qquad f_{*}(Y) = \underline{\mathrm{Hom}}(G, Y),
  \qquad Y \longmapsto \underline{\mathrm{Hom}}(G, Y) .
\]
\marginal{avec noyaux de couples}

Soit $B$ un topos : $\exists\ B'$ topos avec assez de points,
$C \subset \hat{C}$,

\begin{tikzcd}[column sep=small, row sep=small, nodes={font=\scriptsize}]
\simeq B \arrow[d, "f"] & B \arrow[d, "f_{*}", bend left=25] \\
B' & B' \arrow[u, "f^{*}", bend left=25]
\end{tikzcd}

\page{36}
\note{cahier à spirale, papier quadrillé : le feuillet porte une carte des
implications entre classes d'anneaux, reprise et surchargée plusieurs fois. La
partie supérieure est nette et transcrite ci-dessous ; la partie inférieure, où
les mêmes noms sont redistribués sur un second réseau, a ses flèches repassées
au point d'être indéchiffrables — seuls les noms de ses sommets sont donnés}

\begin{tikzcd}[column sep=small, row sep=small, nodes={font=\scriptsize}]
 & \text{anneau} & & \\
\text{anneau local} \arrow[ur] & & \text{anneau réduit} \arrow[ul] & \text{spectre compact} \arrow[l] \\
\text{anneau str. local} \arrow[u] & & \text{anneau intègre} \arrow[u] & \text{pseudocorps} \arrow[ul] \arrow[u] \\
 & & \text{corps} \arrow[u] \arrow[ur] &
\end{tikzcd}

Sommets du second réseau, tels qu'ils se lisent : \emph{spect. connexe},
\emph{\struck{loc. intègre}}, \emph{spectre irréd.}, \emph{réduit},
\emph{spectre compact}, \emph{intègre}, \emph{ps. corps}, \emph{intègre
normal}, \emph{régulier disc.}, \emph{corps} \add{(régulier dim. \ill{})},
\emph{local}, \emph{str. local}, \emph{ultra local}, \emph{loc. irréd.},
\emph{irréd.}, \emph{\ill{} intègre}, \emph{complète}, \emph{\ill{} des anneaux
\ill{} quasi-\ill{}}.

À droite du réseau, la vérification du cas « anneau réduit » :
\[
  \emptyset \xrightarrow{\ \sim\ } I(u) \subset A, \qquad
  \forall f(b)\ \exists\ \struck{e} \in \Pi(u) \text{ tel que } \ldots
\]
\[
  e\,(1-e), \qquad f^{n},\ f, \qquad V(f) \longrightarrow V(f^{n}) ,
\]
\[
  fg = e, \qquad f(1-e) = 0, \qquad f^{n}(1-e) = 0, \qquad f^{n} \overset{?}{=} e .
\]
\note{cette colonne est écrite le feuillet tourné}

En marge droite : $\tilde{R}_{\mathrm{zar}}$, $(A[T]/F_{A}[T])_{F'}$ ;
$X_{i} \to Y$ ; $\mathbb{A}^{n}$, $U$ ; $F \to G$ ; $\Pi \subset A \to A$,
$p \mapsto p^{2} - p$, $\Pi \to e$ \add{(fini ?)}.

\page{37}
\[
  X_{\infty} \subset {}^{e}\hat{R}, \qquad \text{crible de } R,
  \qquad \struck{X_{\infty}^{\mathrm{af}}} .
\]
\[
  X_{\infty}(\mathbb{Z}) =
  \begin{cases}
    \emptyset & \text{si } Z_{\infty} \struck{\ill{}}\ \emptyset \\
    \{e\} & \text{si } Z_{\infty} \neq \emptyset
  \end{cases}
  \qquad
  \underline{\mathrm{Hom}}_{\mathrm{Sch.}\,H \simeq \hat{R}}
    \bigl(\mathbb{Z}, \mathrm{Spec}(\mathbb{Q})\bigr) .
\]
\marginal{ex. g. : $\mathrm{Pro}\,R \to R$, $R \to R$}

\textbf{Sous-topos \struck{ouverts} de $\hat{R} = B$}\quad sous-objet \ill{} de
$e_{B}$, $B' \to B$ (foncteur ponctuel dans $R^{\circ}$) ; crible $C$ de $R$,
$C_{R} \xrightarrow{\ t\ } C$ ; crible \struck{$C$} de $B$ ; sous-structure
distinguée $T'$, $T'_{i} = T_{B}$.
\[
  B' = B/S \xrightarrow{\ \mathrm{can}\ } B, \qquad
  = \hat{C}_{R} \xrightarrow{\ i\ } B = \hat{R}, \qquad
  = \hat{C}_{B}\ \{\varepsilon \xrightarrow{\ \varphi\ } \ill{}\ \text{
    (avec hyp. can.)} \} \longrightarrow \ill{}
\]
soit le « type classifiant » de $T'$, avec
\[
  B_{T'} \xrightarrow{\ \sim\ } \tilde{B} = \hat{B}\ \ill{} , \qquad
  S = i_{!}(\mathbb{1}_{B'}) \ \struck{\ill{}}
\]
\note{la fin de la page est écrite le feuillet tourné d'un quart de tour ; s'y
lit :}
\[
  C_{B} = \mathrm{Id}\ \text{ens.}\ i_{!}
  \;=\; B/S \xrightarrow{\ i\ } \hat{R}
  \;=\; \{\, x \in B \mid A \mid x \in T'(B_{/x}) \,\}
\]
\[
  T' = T_{B} \cdot (T_{i} \to T_{B}) = T_{B/S}, \qquad
  T'(\varepsilon) = \{\, \varepsilon \xrightarrow{\ f\ } B,\ f^{*}(S) \simeq
    e_{\varepsilon} \,\}
\]
avec, plus à gauche, $S \to e_{B}$ un \ill{}morphisme, et
$\chi_{i}^{*}(S) = e_{B'}$.

\page{38}
\note{la page dresse la liste des axiomes sur les anneaux et, en regard, le
topos classifiant de chacun — c'est le cœur de ce cahier}

\begin{enumerate}
  \item[1)] \underline{anneau} (commut. unifère) —
    $X_{\mathrm{ann}} = \hat{R}$, où $R = $ schémas affines de t.f. $/\mathbb{Z}$
    $\simeq$ (algèbres de t.f. sur $\mathbb{Z}$)$^{\circ}$ ;
  \item[2)] \underline{anneau local} ($U \simeq U_{f} \cup U_{1-f}$, cas
    universel $U = A$, $f$ diagonale) —
    $X_{\mathrm{anloc}} = \tilde{R}_{\mathrm{zar}}$, zar $=$ topologie de
    Zariski ;
  \item[4)] \underline{anneau str. local} (local, et pour tout $U$ et sch. rel.
    étale $U'$ sur $U$ avec $U' \to U$ \struck{surjectif} \add{(et schémas
    finis)}, « surjectif » de schémas, \struck{isomorphisme} $U'$ a loc. une
    section sur $U$ ; \ill{} le cas universel,
    $F(T) = aT^{n} + a_{1}T^{n-1} + \cdots + a_{n}$, les $a_{i}$ les $n$
    projections de $A$) —
    $X_{\mathrm{anstrloc}} = \tilde{R}_{\mathrm{et}}$, et $=$ topologie étale ;
    $U = \mathbb{A}^{n}$, $U' = \mathrm{Spec}\ \mathrm{rel} 
    \bigl(\mathcal{O}_{\mathbb{A}^{n}}[T]/(F)\bigr)_{F'}$, $F'$ dérivée de $F$ ;
  \item[3)] \underline{anneau local hensélien} [si $X' \to X$ dans $R$ est
    \ill{} \ill{}, $B' \leftarrow B$ leurs \ill{}, on veut que
    $\underline{\mathrm{Hom}}_{\mathrm{an}}(B', A) \to
    \underline{\mathrm{Hom}}_{\mathrm{an}}(B, A)$ soit épi] —
    $X_{\mathrm{anlochen}} = \tilde{R}_{\mathrm{hen}}$, hen $=$ topologie
    hensélienne : $X_{i} \to X$ \ill{} couvrantes ssi $\forall$ anneau local
    hensélien $A$, $X_{i}(A) \to X(A)$ l'est ;
  \item[5)] \underline{anneau ultralocal} (si $X' \to X$ \add{dans $R$}
    \struck{est} couvrant pour fppf, alors $A(X') \to A(X)$ est épi ; \ill{} et
    dans les cas universels $U = \ill{}$ dans \ill{}. \underline{Question} :
    soit $X' \to \ill{}$ ultralocal ordinaire $A$, $X'(A) \to \ill{}$
    couvrant fppf ? Résultant \ill{} ; toute \ill{} \struck{Rés.} de t.f.
    \ill{} sur $\mathbb{Z}$ admet un \ill{} ultralocal $\mathcal{O}$ qui est
    limite inductive d'algèbres fid. plates de t.f. sur $A$) —
    $X_{\mathrm{anultloc}} = \tilde{R}_{\mathrm{fppf}}$ ; fppf : \ill{} les
    familles $X' \to X$ qu. fini plat et \ill{} fini et plat, $U =
    \mathbb{A}^{n}$, $U' = \mathrm{Spec}\ \struck{\ill{}} 
    \mathcal{O}_{\mathbb{A}^{n}}[T]/(F)$ ;
  \item[6)] \underline{anneau réduit}
    \[
      \mathcal{O}_{A} \xrightarrow{\ \sim\ }
      \varinjlim_{n} \Bigl(\mathrm{Ker}\bigl(A \xrightarrow{\ x \mapsto x^{n}\ }
      A\bigr)\Bigr) \quad [\, I(u) \,]
    \]
    cas universel : $\forall\ \struck{\ill{}}\ V(f^{n}) \to V(f)$,
    $\forall f \in A(U)$, $n \in \mathbb{N}$ —
    $X_{\mathrm{red}} \simeq \tilde{R}_{\mathrm{topred}} \simeq
    \hat{R}_{\mathrm{red}}$ ; topred $=$ la top. sur $R$ où $X_{i} \to Y$ est
    couvrante ssi $\forall\ \mathcal{O}$ réduit, $X_{i}(\mathcal{O}) \to
    Y(\mathcal{O})$ l'est, i.e. s'il existe une section sur $Y_{\mathrm{red}}$ ;
    $R_{\mathrm{red}}$ catégorie des \struck{\ill{}} $R$ réduits (peut être
    aussi définie comme catégorie de fractions de $R$) ;
\end{enumerate}

7) \underline{anneau intègre} : a) $A \amalg A \xrightarrow{\ u\ } D \subset
A \times A$, $x \mapsto (x, 0)$, $y \mapsto (0, y)$, le carré

\begin{tikzcd}[column sep=small, row sep=small, nodes={font=\scriptsize}]
D \arrow[d] \arrow[r, hook] & A \times A \arrow[d, "(x{,}y) \mapsto xy"] \\
e \arrow[r, "0"] & A
\end{tikzcd}

étant cartésien, et $u$ un épi ; b) $\emptyset \xrightarrow{\ \sim\ }
\mathrm{Ker}\bigl(0_{A}, 1_{A} : e \rightrightarrows A\bigr)$. — Topos
classifiant : $X_{\mathrm{int}} \simeq \tilde{R}_{\mathrm{topint}}$ ; topint
$=$ la top. sur $R$ où $X_{i} \to Y$ est couvrante ssi pour toute comp. irréd.
$Y_{\alpha}$ de $Y$ avec str. ind. réduite, il y a section au-dessus.

\marginal{N.B. b) assure que la famille vide \ill{} \ill{} de $R$ \ill{} est
\ill{} pour la top. \ill{}}

\textbf{N.B.}\quad $X_{\mathrm{red}}$ peut être \ill{}, ce qui se voit sur
$X_{\mathrm{red}} \simeq \hat{R}_{\mathrm{red}}$ \ill{} des lim. finies ;
$X_{\mathrm{red}} \simeq \hat{R}_{\mathrm{red}}$ avec $R_{\mathrm{red}}$
\ill{} finies.

Les objets $F \in \tilde{R}_{\mathrm{topint}}$ sont les $F \in \hat{R}$ tels
que $\forall X \in \mathrm{Ob}\,R$ et $(X_{i})_{i \in I}$ famille des comp.
irréd.,
\[
  F(X) \longrightarrow \prod_{i} F(X_{i}) \rightrightarrows
  \prod_{i,j} F(X_{i} \times_{X} X_{j})
\]
soit exacte.

\page{39}
8) \underline{Anneau irréductible} : $A_{\mathrm{red}}
\overset{\text{déf}}{=} A \big/ \varinjlim_{n} \bigl(I_{n} =
\mathrm{Ker}(A \xrightarrow{\ x \mapsto x^{n}\ } A)\bigr)$ \add{nilradical
$N$} est intègre, i.e. a) $A \times N \amalg N \times A \xrightarrow{\ v\ }
D_{N} \hookrightarrow A \times A$, le carré

\begin{tikzcd}[column sep=small, row sep=small, nodes={font=\scriptsize}]
D_{N} \arrow[d] \arrow[r, hook] & A \times A \arrow[d, "(x{,}y) \mapsto xy"] \\
N \arrow[r, hook] & A
\end{tikzcd}

étant cartésien, et $v$ épi ; b) $V(1_{A}) \xrightarrow{\ \simeq\ }
\emptyset_{\varepsilon}$, i.e. \struck{\ill{}}
\[
  \begin{cases}
    \forall f, g \in \struck{A}(U), \quad
    \varinjlim_{n} V(f^{n}g^{n}) \Longleftarrow
    \varinjlim_{n} V(f^{n}) \vee \varinjlim_{n} V(g^{n}) \\
    V(1_{A}) \simeq \emptyset_{\varepsilon}
  \end{cases}
\]

\marginal{on obtient un exemple d'un $X_{\mathrm{irr}}$ \struck{dans le
\ill{}} couvrant du couvrant $X_{\mathrm{ann}}$}

\[
  X_{\mathrm{irr}} \simeq \tilde{R}_{\mathrm{topirr}}
\]
topirr $=$ la topologie où une famille $X_{i} \to X$ est couvrante ssi
\struck{\ill{}} $\forall$ comp. irréductible $Y$ de $X$, \ill{} \ill{} ;
th. : l'un des $X_{i}$ a une section sur $Y_{\mathrm{red}}$.

\textbf{N.B.}\quad Cette topologie n'est pas quasi-compacte : il y a des
recouvrements (de \struck{l'\ill{}} tout schéma \struck{ayant une}
\add{d'\ill{} comp. connexes} \struck{\ill{}} irréductibles) qui n'admettent
pas de ss-recouvrement fini, \ill{} \ill{} finies en produisant
$R_{\mathrm{ann}}$ ($R_{\mathrm{ann}}$ ss.-cat. des \ill{} conn.).

\textbf{N.B.}\quad Les $F \in \struck{\mathrm{Ob}}\ X_{\mathrm{irr}}$ \ill{}
$\hat{R}$ \ill{} de $\tilde{R}_{\mathrm{topirr}}$ \ill{}, d'un
\struck{sous-}topos \ill{} du topos \ill{}

\page{40}
\note{une seule ligne sur ce feuillet, portée le cahier retourné}
\ill{} $R$ à $\varepsilon$ ont $X_{\emptyset} =$ sous-topos
\uncertain{fermé} \ill{} de $X_{\mathrm{ann}}$, correspondant à l'anneau
\uncertain{unité} et \ill{}

\end{document}
